\documentclass[10pt,a4paper]{amsart}
\usepackage[margin=19mm]{geometry}
\usepackage{amsmath,amssymb,mathtools}
\usepackage[scr=rsfs,cal=euler]{mathalfa}
\usepackage{xfrac}
\usepackage[unicode,hidelinks,bookmarks=false]{hyperref}
\newcommand{\ie}{i.e.\ }
\newcommand{\cf}{cf.\ }
\newcommand{\resp}{respectively\ }
\newcommand{\Subsection}{\subsection}
\providecommand{\nobreakdash}{\nobreak\hbox{-}}
\setlength{\emergencystretch}{2em}
\multlinegap0pt
\usepackage{amssymb}
\usepackage{enumerate}
\usepackage{mathtools}
\usepackage{xcolor}
\usepackage{float}
\newtheorem{lemma}{Lemma}
\newcommand{\GG}{\mathcal{G}}
\title{The Gill--Guillot commuting graph \\for sporadic and related groups\vspace*{-0.4cm}}
\author{David A. Craven, Coen del Valle, Chris Parker}
\date{Source: arXiv:2602.02032v2 (CC BY 4.0)}
\begin{document}
\maketitle

\noindent\emph{Source and licence.} The Gill--Guillot commuting graph \\for sporadic and related groups\vspace*{-0.4cm}. Version arXiv:2602.02032v2 (CC BY 4.0), distributed by arXiv under the Creative Commons Attribution 4.0 International licence, http://creativecommons.org/licenses/by/4.0/, as recorded in the arXiv licence metadata for this exact version and shown on the abstract page https://arxiv.org/abs/2602.02032v2. Full licence notice: \texttt{licenses/arxiv-2602.02032-CC-BY-4.0.txt}.\par\smallskip

\noindent\textit{Editorial scope.} Counted unit: the article's lemma lem:covermulticlass (source0.tex lines 440--443). The article's complete original proof of it is delivered with it (source0.tex lines 444--450, one \texttt{proof} environment of 7 lines). No mathematics is written, repaired or re-derived: the statement and the proof are the article's own lines, in the article's own notation, and no retained line is substituted or re-flowed.  No further source-local context is needed: every cross-reference of every retained line resolves inside the two delivered spans.  Nothing was omitted from any retained span, so the package carries no omission record.  No retained line contains a citation, so the wrapper carries no bibliography and no bibliography entry.  No retained line contains a figure environment, an image-inclusion command or drawing code, so no image is delivered and the package declares no asset dependency.  Editorial formatting only, in addition: an AMS \texttt{amsart} preamble that restates the article's own declarations and macros used by the retained lines, namely the environments \texttt{lemma} on separate counters, together with the standard packages those lines need.  The retained environments print in this document as Lemma 1 in order of appearance, whereas the article prints the same items with the numbering of its own full document.  The complete native source of the article as distributed in the arXiv e-print for this exact version is bundled unchanged as \texttt{sources/193821/source0.tex}.
\begin{lemma}\label{lem:covermulticlass}

Suppose that $G$ is a group and $Z \le Z(G)$ is a $2$-group. Set $\widetilde{G}=G/Z$ and let $\varphi$ be the quotient map from $G$ to $\widetilde{G}$. Suppose that $\widetilde{\mathcal{C}}$ is a conjugacy class of non-central involutions in $\widetilde G$ such that $ \varphi^{-1}(\widetilde{\mathcal{C}})= \bigsqcup_{j=1}^m\mathcal{C}_j $ is a union of $m>1$ conjugacy classes of involutions in $G$. Assume that $\GG(\widetilde{\mathcal{C}})$ is connected. If $n_G(\mathcal{C}_1,\mathcal{C}_1,\mathcal{C}_j)=0$ for all $1< j\leq m$, then $\GG(\mathcal{C}_1)$ is connected and $\GG(\mathcal{C}_j)$ is edgeless for all $1<j\leq m$.
\end{lemma}

\begin{proof} Suppose that $\{Zx,Zy\}$ is an edge in $\GG(\widetilde{\mathcal C})$. Let $u\in \mathcal C_1 \cap Zx$ and $v \in \mathcal C_1 \cap Zy$. Then $Z xy\in \widetilde {\mathcal{C}}$, hence $uv \in \bigsqcup_{j=1}^m\mathcal{C}_j$. As $n_G(\mathcal{C}_1,\mathcal{C}_1,\mathcal{C}_j)=0$ for all $1<j\le m$, we have $uv \in \mathcal C_1$. In particular, $uv$ is an involution and consequently $u$ and $v$ commute. Hence $u$ and $v$ are adjacent in $\GG(\mathcal{C}_1)$. It follows that paths in $\GG(\widetilde{\mathcal{C}})$ lift to paths in $\GG(\mathcal{C}_1)$, and so
$\GG(\mathcal {C}_1)$ is connected.

%Let $x,y\in\mathcal{C}_1$, and let $\varphi(x):=v_1,\dots,v_k=:\varphi(y)$ be a path in $\GG(\widetilde{\mathcal{C}})$. Set $u_1=x$ and $u_k=y$, and for each $i\in\{2,\dots,k-1\}$ set $u_i\in\varphi^{-1}(v_i)\cap\mathcal{C}_1$. Then for each $i\in\{1,\dots,k-1\}$ we have that $\varphi(u_iu_{i+1})=v_iv_{i+1}\in\widetilde{\mathcal{C}}$, whence $u_iu_{i+1}\in\bigsqcup_{j=1}^m\mathcal{C}_i$. Since $n_G(\mathcal{C}_1,\mathcal{C}_1,\mathcal{C}_j)=0$ for all $1< j\leq m$, we deduce that $u_iu_{i+1}\in\mathcal{C}_1$ and so we have a path from $x$ to $y$ as desired.

For the second claim suppose for contradiction that $j>1$ and that there exist $w,r,s\in \mathcal{C}_j$ satisfying $rs=w$. Take $z\in Z$ such that $rz\in\mathcal{C}_1$. Then $1=(rz)^2=r^2z^2=z^2$ and so $z$ is an involution. As $r$ and $s$ are conjugate and $z$ is central, we have $sz\in \mathcal C_1$. Now  $(rz)(sz)=rsz^2=rs=w$, contradicting the assumption that $n_G(\mathcal{C}_1,\mathcal{C}_1,\mathcal{C}_j)=0$, hence the result.
\end{proof}

\end{document}
